We simulated the sandpile model for the configurations listed in Table \ref{tab:model_configurations}. A burn-in procedure was used to reach the stationary critical state. Monitoring $\expval{z}(\tau)$ shows a clear approach to a plateau (Fig. \ref{fig:burn-in-comparison-zero}). Testing different initializations indicates that starting from $z_0=z_c$ reduces the burn-in time for most systems (Fig. \ref{fig:burn-in-comparison-init}, Table \ref{tab:burn_in_times}).

In the critical state, the PDFs of avalanche lifetime $t$, size $s$, and linear extension $l$ show approximate power-law scaling over an intermediate range (example fits in Fig. \ref{fig:pdf-plot-example}). The resulting PDF exponents $\alpha$, $\tau$, and $\lambda$ are summarized for open nonconservative, open conservative, and closed nonconservative driving in Tables \ref{tab:scaling_exp_open_nonconservative}--\ref{tab:scaling_exp_closed_nonconservative}. Conditional scaling exponents from the expectation value fits (Fig. \ref{fig:cond-plot-example}) are listed in Tables \ref{tab:scaling_exp_open_nonconservative_gamma}--\ref{tab:scaling_exp_closed_nonconservative_gamma}. Overall, the exponent trends with dimension and with the driving scheme are consistent with the literature within uncertainties, with the largest uncertainties in $d=4$ where the scaling range is shortest. For follow-up work, larger system sizes and more measurements in $d=3$ and $d=4$ should improve the scaling windows, especially for observables involving $l$.
